% Libro 1 · Mundo 4 · Proporcionalidad y porcentajes (app: pu-proporcion)
\mundo{Proporcionalidad y porcentajes}{pu-proporcion}{Una receta para cuatro que se convierte en una para diez, un
mapa que encoge una ciudad, dos obreros que terminan antes que uno, un descuento sobre un descuento. Todo eso es
proporcionalidad: comparar cantidades que cambian juntas. Es, con diferencia, el tema que más preguntas aporta al
examen.}{%
\item \textbf{4.1} Razones, repartos y escalas
\item \textbf{4.2} Regla de tres simple y compuesta
\item \textbf{4.3} Velocidad y trabajo conjunto
\item \textbf{4.4} Porcentajes: las tres preguntas
\item \textbf{4.5} Aumentos, descuentos e IVA}

% ------------------------------------------------------------------
\leccion{pu-pr-1}{Razones, repartos y escalas}{Juegos de la limonada y la balanza de proporciones.}

\subsection{La idea: partes iguales}
Decir que hombres y mujeres están en razón $3 : 5$ es decir que el curso se puede partir en \clave{grupos iguales}: por
cada 3 hombres hay 5 mujeres. No dice cuántos hay; dice cómo se comparan. La herramienta para verlo es la
\clave{barra de partes}:

\figura{razon}{Razón 3 : 5 en barras: conocer una parte es conocerlo todo.}

Cuando se sabe cuánto vale \emph{una} parte, se sabe todo. Este es el corazón de los repartos y las escalas.

\subsubsection{Proporción}
Dos razones iguales forman una \clave{proporción}: $\tfrac ab = \tfrac cd$. Su propiedad clave es que los productos
cruzados son iguales: $ad = bc$. Si falta un término, se despeja: en $\tfrac{3}{5} = \tfrac{x}{40}$, $x = \tfrac{3 \cdot
40}{5} = 24$.

\subsubsection{Reparto proporcional}
Repartir $N$ en razón $a : b : c$ es partir $N$ en $a + b + c$ partes iguales y dar $a$, $b$ y $c$ de ellas. Si el
reparto es \emph{inversamente} proporcional (por ejemplo, más ayuda a quien menos tiene), se reparte en razón
$\tfrac1a : \tfrac1b : \tfrac1c$, que conviene convertir a enteros multiplicando por un múltiplo común.

\subsubsection{Escalas}
Una escala $1 : 50\,000$ dice que 1 unidad en el mapa son 50\,000 de esas mismas unidades en la realidad: 1 cm son
50\,000 cm, es decir, 500 m. Para áreas, la escala se eleva al cuadrado: 1 cm$^2$ del mapa son $50\,000^2$ cm$^2$
reales.

\begin{regla}
\begin{itemize}
  \item Razón $a : b$ y total $N$: cada parte vale $\tfrac{N}{a + b}$.
  \item Razón $a : b$ y diferencia $D$: cada parte vale $\tfrac{D}{b - a}$.
  \item $\tfrac ab = \tfrac cd \iff ad = bc$.
  \item Reparto inverso a $a, b, c$: en razón $\tfrac1a : \tfrac1b : \tfrac1c$.
  \item Escala $1 : k$: longitudes por $k$, áreas por $k^2$, volúmenes por $k^3$.
\end{itemize}
\end{regla}

\subsection{Ejemplos resueltos}
\begin{ejemplo}[Reparto en tres]
Tres socios invierten 4, 6 y 10 millones y reparten una ganancia de 15 millones en proporción a lo invertido. ¿Cuánto
recibe cada uno?
\solucion La razón $4 : 6 : 10$ se simplifica a $2 : 3 : 5$, con $10$ partes. Cada parte vale $1.5$ millones: reciben
$3$, $4.5$ y $7.5$ millones.
\end{ejemplo}

\begin{ejemplo}[Reparto inverso]
Se reparten 260 cuadernos entre tres escuelas en forma inversamente proporcional a su número de computadores: 2, 3 y 4.
¿Cuántos recibe la que tiene 3?
\solucion Razón $\tfrac12 : \tfrac13 : \tfrac14$. Multiplicando por 12: $6 : 4 : 3$, con 13 partes de $260 \div 13 =
20$. La escuela de 3 computadores recibe $4 \cdot 20 = 80$.
\end{ejemplo}

\begin{ejemplo}[Mezclas]
Una pintura verde se hace con azul y amarillo en razón $2 : 3$. Si hay 8 litros de azul, ¿cuánto amarillo falta para
usarlos todos, y cuánta pintura verde resulta?
\solucion 2 partes son 8 L, una parte 4 L. Amarillo: $3 \cdot 4 = 12$ L. Verde: $5 \cdot 4 = 20$ L.
\end{ejemplo}

\begin{ejemplo}[Escala y área]
En un plano a escala $1 : 200$, un salón mide $4 \text{ cm} \times 3 \text{ cm}$. ¿Cuál es su área real?
\solucion Las medidas reales son $8$ m y $6$ m: $48$ m$^2$. (Con la escala al cuadrado: $12 \text{ cm}^2 \cdot 40\,000
= 480\,000$ cm$^2 = 48$ m$^2$.)
\end{ejemplo}

\begin{trampa}
En una razón $3 : 5$, los hombres no son $\tfrac35$ del curso: son $\tfrac38$ (3 de cada 8 personas). La razón compara
las partes entre sí; la fracción compara una parte con el total.
\end{trampa}

\begin{atajo}
Antes de repartir, simplifica la razón: $4 : 6 : 10$ es $2 : 3 : 5$. Menos partes, cuentas más cortas.
\end{atajo}

\begin{reto}
En un curso, la razón entre hombres y mujeres es $3 : 5$. Si hay 12 mujeres más que hombres, ¿cuántos estudiantes hay?
\opciones{32}{40}{48}{56}
\end{reto}

\begin{retoresuelto}
Las mujeres tienen 2 partes más que los hombres ($5 - 3$), y esas 2 partes son 12 personas: cada parte vale 6. El curso
tiene $3 + 5 = 8$ partes: $8 \cdot 6 = 48$ estudiantes. \textbf{Opción C.}
\end{retoresuelto}

\taller
\piso{Calienta}{partes}
\begin{ejercicios}
\item Reparte 72 en razón $4 : 5$. \sol{Partes de 8: 32 y 40.}
\item Halla $x$: $\tfrac{x}{12} = \tfrac{15}{20}$. \sol{$x = 9$.}
\item En un mapa $1 : 25\,000$, ¿cuántos km son 6 cm? \sol{$150\,000$ cm $= 1.5$ km.}
\end{ejercicios}
\piso{Entrena}{problemas}
\begin{ejercicios}
\item Las edades de dos hermanos están en razón $2 : 7$ y suman 36. ¿Cuántos años se llevan? \sol{Partes de 4: 8 y 28. Se llevan 20.}
\item Se reparten 330 000 pesos entre tres niños en forma inversamente proporcional a sus edades: 2, 5 y 10 años. ¿Cuánto recibe el menor? \sol{$\tfrac12 : \tfrac15 : \tfrac1{10} = 5 : 2 : 1$; 8 partes de 41\,250. El de 2 años recibe 206\,250.}
\item Un cubo a escala $1 : 10$ tiene volumen de 2 cm$^3$. ¿Cuál es el volumen del cubo real? \sol{$2 \cdot 10^3 = 2\,000$ cm$^3$.}
\end{ejercicios}
\piso{Como en el examen}{opción múltiple}
\begin{ejercicios}
\item En una bolsa hay canicas rojas y azules en razón $4 : 7$. Si hay 21 azules más que rojas, el total de canicas es
\opciones{33}{55}{77}{88}
\sol{C. Tres partes son 21: la parte es 7 y el total $11 \cdot 7 = 77$.}
\item La razón entre $a$ y $b$ es $2 : 3$ y la razón entre $b$ y $c$ es $4 : 5$. La razón entre $a$ y $c$ es
\opciones{$8 : 15$}{$2 : 5$}{$6 : 5$}{$8 : 5$}
\sol{A. Se iguala $b$: $a : b = 8 : 12$, $b : c = 12 : 15$.}
\item En un plano a escala $1 : 500$, un terreno ocupa 20 cm$^2$. Su área real es
\opciones{1\,000 m$^2$}{500 m$^2$}{100 m$^2$}{10\,000 m$^2$}
\sol{B. $20 \cdot 500^2$ cm$^2 = 5\,000\,000$ cm$^2 = 500$ m$^2$.}
\end{ejercicios}

\bitacora

% ------------------------------------------------------------------
\leccion{pu-pr-2}{Regla de tres simple y compuesta}{Juegos de las magnitudes que suben juntas y la balanza inversa.}

\subsection{La idea: ¿suben juntas o una sube y la otra baja?}
Dos magnitudes son \clave{directamente proporcionales} si al duplicar una se duplica la otra: kilos y precio, horas y
salario. Lo que se conserva es el \emph{cociente} (el precio por kilo). Son \clave{inversamente proporcionales} si al
duplicar una la otra se reduce a la mitad: obreros y días para una misma obra, velocidad y tiempo para un mismo camino.
Lo que se conserva es el \emph{producto} (el trabajo total).

\figura{dirinv}{Directa: una recta que pasa por el origen. Inversa: una curva que baja.}

\subsubsection{La regla de tres compuesta como «trabajo total»}
Cuando intervienen varias magnitudes (obreros, horas por día, días, metros de obra), la forma más segura es calcular una
cantidad que se conserva: el \clave{trabajo} medido en obrero-horas. Si 5 obreros trabajando 8 horas diarias hacen
una obra en 12 días, la obra «cuesta» $5 \cdot 8 \cdot 12 = 480$ obrero-horas. Con esa cifra se responde cualquier otra
pregunta sobre la misma obra.

\begin{regla}
\begin{itemize}
  \item Directa: $\dfrac{y_1}{x_1} = \dfrac{y_2}{x_2}$. Inversa: $x_1 y_1 = x_2 y_2$.
  \item Compuesta: $\dfrac{\text{trabajadores} \times \text{horas} \times \text{días}}{\text{obra}}$ es constante (si
  la eficiencia es la misma).
  \item Para decidir si es directa o inversa, pregúntate: \emph{si hubiera más de esto, ¿haría falta más o menos de lo
  otro?}
\end{itemize}
\end{regla}

\subsection{Ejemplos resueltos}
\begin{ejemplo}[Directa]
Si 3 kg de queso cuestan \$51\,000, ¿cuánto cuestan 5 kg?
\solucion Un kilo vale $17\,000$; cinco kilos, $85\,000$.
\end{ejemplo}

\begin{ejemplo}[Inversa]
Un viaje dura 6 horas a 80 km/h. ¿Cuánto dura a 120 km/h?
\solucion La distancia (producto) es 480 km. A 120 km/h: $480 \div 120 = 4$ horas.
\end{ejemplo}

\begin{ejemplo}[Compuesta]
Ocho máquinas, trabajando 6 horas diarias, fabrican 1\,200 piezas en 5 días. ¿Cuántos días necesitan 10 máquinas,
trabajando 8 horas diarias, para fabricar 2\,400 piezas?
\solucion Trabajo por pieza: $\tfrac{8 \cdot 6 \cdot 5}{1\,200} = \tfrac{240}{1\,200} = 0.2$ máquina-horas. Para 2\,400
piezas hacen falta $480$ máquina-horas. Con $10 \cdot 8 = 80$ máquina-horas por día: $480 \div 80 = 6$ días.
\end{ejemplo}

\begin{ejemplo}[Provisiones]
Un campamento de 30 personas tiene comida para 12 días. Al cabo de 4 días llegan 10 personas más. ¿Para cuántos días
más alcanza la comida?
\solucion Comida total: $30 \cdot 12 = 360$ raciones-día. En 4 días se gastan 120 y quedan 240. Para 40 personas:
$240 \div 40 = 6$ días.
\end{ejemplo}

\begin{trampa}
Aplicar regla de tres directa a magnitudes inversas: «si 4 obreros tardan 6 días, 8 obreros tardan 12 días». El sentido
común lo descarta: más obreros, menos días.
\end{trampa}

\begin{atajo}
En la compuesta, no armes tablas de flechas: calcula el «trabajo total» y divide. Es una sola idea que sirve para todos
los casos.
\end{atajo}

\begin{reto}
Seis grifos iguales llenan un tanque en 8 horas. ¿Cuántas horas tardan 4 de esos grifos en llenar dos tanques iguales
al primero?
\opciones{6}{12}{24}{32}
\end{reto}

\begin{retoresuelto}
Un tanque cuesta $6 \cdot 8 = 48$ grifo-horas. Dos tanques: 96 grifo-horas. Con 4 grifos: $96 \div 4 = 24$ horas.
\textbf{Opción C.} La B ($48 \div 4 = 12$) olvida el segundo tanque.
\end{retoresuelto}

\taller
\piso{Calienta}{directa o inversa}
\begin{ejercicios}
\item 5 cuadernos cuestan \$17\,500. ¿Cuánto cuestan 8? \sol{\$28\,000.}
\item 12 obreros hacen una casa en 30 días. ¿Cuántos días tardan 18? \sol{$360 / 18 = 20$ días.}
\item Un grifo llena un tanque en 4 horas con un caudal de 15 L/min. ¿Cuánto tardaría con 20 L/min? \sol{3 horas.}
\end{ejercicios}
\piso{Entrena}{compuesta}
\begin{ejercicios}
\item 6 pintores pintan 4 casas en 10 días. ¿Cuántos pintores se necesitan para pintar 6 casas en 9 días? \sol{Por casa: $6 \cdot 10 / 4 = 15$ pintor-días. Seis casas: 90. En 9 días: 10 pintores.}
\item Una familia de 4 gasta 120 kg de arroz en 60 días. ¿Cuánto gasta una de 6 en 30 días? \sol{Por persona-día: $\tfrac{120}{240} = 0.5$ kg. $6 \cdot 30 \cdot 0.5 = 90$ kg.}
\item Un barco tiene víveres para 25 días con 40 tripulantes. ¿Para cuántos días alcanzan si se embarcan 10 más? \sol{$1\,000 / 50 = 20$ días.}
\end{ejercicios}
\piso{Como en el examen}{opción múltiple}
\begin{ejercicios}
\item Si 3 gallinas ponen 3 huevos en 3 días, ¿cuántos huevos ponen 12 gallinas en 12 días?
\opciones{12}{24}{36}{48}
\sol{D. Una gallina pone 1 huevo cada 3 días: en 12 días, 4. Doce gallinas: 48.}
\item 20 obreros terminan una obra en 15 días. Después de 5 días se retiran 5 obreros. ¿Cuántos días tardan en total?
\opciones{$18\tfrac13$}{$20$}{$13\tfrac13$}{$18$}
\sol{A. Obra: 300 obrero-días. En 5 días se hacen 100; faltan 200 con 15 obreros: $13\tfrac13$ días más. Total $18\tfrac13$.}
\item Un automóvil gasta 8 galones en 320 km. ¿Cuántos galones gasta en 500 km?
\opciones{10}{12.5}{12}{15}
\sol{B. 40 km por galón: $500 / 40 = 12.5$.}
\end{ejercicios}

\bitacora

% ------------------------------------------------------------------
\leccion{pu-pr-3}{Velocidad y trabajo conjunto}{Juegos de los dos móviles en la pista y los grifos que llenan juntos.}

\subsection{La idea: velocidades que se suman o se restan}
La relación de siempre es $d = v \cdot t$. Lo nuevo aparece con dos móviles: si se \clave{acercan} de frente, la distancia
entre ellos se cierra con la \emph{suma} de sus velocidades; si uno \clave{persigue} al otro en la misma dirección, la
distancia se cierra con la \emph{diferencia}.

\figura{encuentro}{Al acercarse de frente, las velocidades se suman.}

\subsubsection{Trabajo conjunto: se suman las rapideces, no los tiempos}
Si un grifo llena un tanque en 3 horas, llena $\tfrac13$ de tanque por hora. Si otro lo llena en 6, llena $\tfrac16$ por
hora. Juntos llenan $\tfrac13 + \tfrac16 = \tfrac12$ por hora: tardan 2 horas. Lo que se suma es la \clave{rapidez}
(fracción de trabajo por unidad de tiempo); el tiempo conjunto es el inverso de esa suma. Un desagüe que vacía resta su
rapidez.

\begin{regla}
\begin{itemize}
  \item $d = v t$. Encuentro de frente: $t = \dfrac{d}{v_1 + v_2}$. Persecución: $t = \dfrac{d}{v_1 - v_2}$.
  \item Si $A$ hace un trabajo en $a$ horas y $B$ en $b$ horas, juntos tardan
  $\dfrac{1}{\frac1a + \frac1b} = \dfrac{ab}{a + b}$.
  \item Rapidez promedio $=$ distancia total $\div$ tiempo total (no el promedio de las rapideces).
\end{itemize}
\end{regla}

\subsection{Ejemplos resueltos}
\begin{ejemplo}[Persecución]
Un ciclista sale a 15 km/h. Dos horas después sale tras él una moto a 45 km/h. ¿Cuánto tarda la moto en alcanzarlo?
\solucion Cuando sale la moto, el ciclista lleva 30 km. La moto descuenta $45 - 15 = 30$ km cada hora: lo alcanza en
1 hora.
\end{ejemplo}

\begin{ejemplo}[Tres que trabajan]
Ana pinta un muro en 4 horas, Beto en 6 y Caro en 12. ¿Cuánto tardan los tres juntos?
\solucion Rapideces: $\tfrac14 + \tfrac16 + \tfrac{1}{12} = \tfrac{3 + 2 + 1}{12} = \tfrac12$ de muro por hora: 2 horas.
\end{ejemplo}

\begin{ejemplo}[Grifo y desagüe]
Un grifo llena una piscina en 6 horas y el desagüe la vacía en 10. Si se abren los dos con la piscina vacía, ¿cuánto
tarda en llenarse?
\solucion $\tfrac16 - \tfrac1{10} = \tfrac{5 - 3}{30} = \tfrac{1}{15}$ por hora: 15 horas.
\end{ejemplo}

\begin{ejemplo}[Ida y vuelta]
Un atleta sube una montaña a 4 km/h y baja por el mismo camino a 12 km/h. ¿Cuál es su rapidez promedio?
\solucion Supón 12 km de camino: sube en 3 h y baja en 1 h. Recorre 24 km en 4 h: 6 km/h (no 8).
\end{ejemplo}

\begin{trampa}
Promediar tiempos en el trabajo conjunto: «4 horas y 6 horas, juntos 5 horas». Juntos siempre tardan \emph{menos} que el
más rápido solo.
\end{trampa}

\begin{atajo}
Para dos personas, $\dfrac{ab}{a + b}$ da el tiempo conjunto de inmediato: 4 y 12 horas $\to \tfrac{48}{16} = 3$ horas.
\end{atajo}

\begin{reto}
Dos ciudades están a 360 km. A la misma hora sale un bus de cada una, hacia la otra: uno a 70 km/h y el otro a 50 km/h.
¿En cuántas horas se encuentran?
\opciones{2}{3}{3.6}{5}
\end{reto}

\begin{retoresuelto}
Se acercan $70 + 50 = 120$ km cada hora; los 360 km se cierran en $360 \div 120 = 3$ horas. \textbf{Opción B.} La D ($360 \div 70 \approx 5$) solo tiene en cuenta el bus rápido; usar la diferencia de velocidades, que es para persecuciones, daría 18 horas.
\end{retoresuelto}

\taller
\piso{Calienta}{cálculos directos}
\begin{ejercicios}
\item ¿Cuánto tarda en recorrer 210 km un carro a 70 km/h? \sol{3 h.}
\item Juan hace una tarea en 3 horas y Pedro en 6. ¿Juntos? \sol{$\tfrac{18}{9} = 2$ h.}
\item Dos personas a 600 m caminan una hacia la otra a 1.5 m/s y 2.5 m/s. ¿En cuánto se encuentran? \sol{$600 / 4 = 150$ s.}
\end{ejercicios}
\piso{Entrena}{problemas}
\begin{ejercicios}
\item Un tren de 200 m a 20 m/s cruza un puente de 400 m. ¿Cuánto tarda en cruzarlo por completo? \sol{Debe recorrer $400 + 200 = 600$ m: 30 s.}
\item Dos grifos juntos llenan un tanque en 4 horas. Uno solo lo llena en 6. ¿Cuánto tarda el otro solo? \sol{$\tfrac14 - \tfrac16 = \tfrac1{12}$: 12 horas.}
\item Un nadador avanza 3 km/h en aguas quietas y el río corre a 1 km/h. ¿Cuánto tarda en ir 4 km río arriba y volver? \sol{Sube a 2 km/h (2 h) y baja a 4 km/h (1 h): 3 horas.}
\end{ejercicios}
\piso{Como en el examen}{opción múltiple}
\begin{ejercicios}
\item Una llave llena un tanque en 2 horas y otra en 3 horas. Juntas lo llenan en
\opciones{2.5 h}{1.2 h}{1.5 h}{5 h}
\sol{B. $\tfrac{6}{5} = 1.2$ h.}
\item Un carro va de A a B a 60 km/h y regresa a 40 km/h. Su rapidez promedio es
\opciones{50 km/h}{48 km/h}{45 km/h}{52 km/h}
\sol{B. Con 120 km de distancia: 2 h + 3 h para 240 km.}
\item Dos corredores dan vueltas a una pista circular de 400 m en el mismo sentido, a 5 m/s y 3 m/s, saliendo juntos. ¿Cada cuánto el rápido le saca una vuelta al lento?
\opciones{50 s}{80 s}{200 s}{133 s}
\sol{C. Le gana 2 m por segundo; una vuelta son 400 m: 200 s.}
\end{ejercicios}

\bitacora

% ------------------------------------------------------------------
\leccion{pu-pr-4}{Porcentajes}{Juegos de pintar personas y de la barra del cien por ciento.}

\subsection{La idea: tres preguntas distintas}
Con porcentajes se hacen tres preguntas, y la mitad de los errores vienen de confundirlas:

\begin{center}
\begin{tabular}{@{}lll@{}}
\toprule
\textbf{Pregunta} & \textbf{Ejemplo} & \textbf{Cuenta}\\\midrule
¿Cuánto es el $p\,\%$ de $N$? & el $30\,\%$ de 80 & $0.3 \cdot 80 = 24$\\
¿Qué porcentaje es $A$ de $B$? & ¿qué \% es 18 de 72? & $\tfrac{18}{72} = 25\,\%$\\
¿De qué número $A$ es el $p\,\%$? & 45 es el $15\,\%$ de… & $45 \div 0.15 = 300$\\
\bottomrule
\end{tabular}
\end{center}

La tercera es la más traicionera, y la que se resuelve mejor con la \clave{barra del cien por ciento}: si 45 es el
$15\,\%$, el $5\,\%$ es 15 y el $100\,\%$ es $15 \cdot 20 = 300$.

\subsubsection{Porcentaje de un porcentaje}
El $20\,\%$ del $30\,\%$ de algo es $0.2 \cdot 0.3 = 0.06$, es decir, el $6\,\%$. Los porcentajes encadenados se
multiplican.

\subsubsection{Puntos porcentuales}
Si el desempleo pasa del $10\,\%$ al $12\,\%$, subió 2 \emph{puntos porcentuales}, pero creció un $20\,\%$
(porque 2 es el $20\,\%$ de 10). Son dos formas de medir el mismo cambio, y el examen las distingue.

\figura{barracien}{La barra del cien por ciento.}

\begin{regla}
\begin{itemize}
  \item Parte $= \tfrac{p}{100}\cdot$ total. Porcentaje $= \tfrac{\text{parte}}{\text{total}}\cdot 100$.
  Total $= \text{parte} \div \tfrac{p}{100}$.
  \item El $a\,\%$ del $b\,\%$ es el $\tfrac{ab}{100}\,\%$.
  \item La «base» del porcentaje es siempre la cantidad que sigue a la palabra «de».
\end{itemize}
\end{regla}

\subsection{Ejemplos resueltos}
\begin{ejemplo}[¿Qué parte del total?]
En un colegio hay 480 estudiantes; el $45\,\%$ son mujeres y el $25\,\%$ de las mujeres practica un deporte. ¿Cuántas
mujeres deportistas hay?
\solucion $0.25 \cdot 0.45 = 0.1125$; el $11.25\,\%$ de 480 es $54$. O por pasos: 216 mujeres, la cuarta parte 54.
\end{ejemplo}

\begin{ejemplo}[El total desde una parte]
Pagué \$63\,000, que era el $70\,\%$ del precio original. ¿Cuál era?
\solucion El $10\,\%$ es $9\,000$; el $100\,\%$, $90\,000$.
\end{ejemplo}

\begin{ejemplo}[Porcentaje de cambio]
Una planta pasó de 40 cm a 52 cm. ¿En qué porcentaje creció?
\solucion Creció 12 cm sobre 40 cm: $\tfrac{12}{40} = 30\,\%$. La base es el valor \emph{inicial}.
\end{ejemplo}

\begin{ejemplo}[Mezcla de porcentajes]
Se mezclan 30 litros de jugo con $20\,\%$ de azúcar y 20 litros con $45\,\%$. ¿Qué porcentaje de azúcar tiene la
mezcla?
\solucion Azúcar: $6 + 9 = 15$ litros en 50 litros: $30\,\%$. (Es un promedio ponderado: más cerca de $20\,\%$ porque
hay más de ese jugo.)
\end{ejemplo}

\begin{trampa}
Calcular un porcentaje de cambio sobre el valor final. De 40 a 52 es un aumento del $30\,\%$; de 52 a 40 es una
disminución de $\tfrac{12}{52} \approx 23\,\%$, no del $30\,\%$.
\end{trampa}

\begin{reto}
En unas elecciones, el candidato A obtuvo el $40\,\%$ de los votos, el candidato B obtuvo $\tfrac38$ y los 9\,000 votos
restantes fueron en blanco. ¿Cuántos votos hubo?
\opciones{36\,000}{40\,000}{45\,000}{24\,000}
\end{reto}

\begin{retoresuelto}
$40\,\% = \tfrac25$. Entre A y B: $\tfrac25 + \tfrac38 = \tfrac{16 + 15}{40} = \tfrac{31}{40}$. En blanco: $\tfrac{9}{40}$,
que son 9\,000 votos. Un cuarentavo son 1\,000 y el total 40\,000. \textbf{Opción B.}
\end{retoresuelto}

\taller
\piso{Calienta}{las tres preguntas}
\begin{ejercicios}
\item ¿Cuánto es el $35\,\%$ de 260? \sol{91.}
\item ¿Qué porcentaje es 21 de 60? \sol{$35\,\%$.}
\item 56 es el $80\,\%$ de un número. ¿Cuál? \sol{70.}
\end{ejercicios}
\piso{Entrena}{problemas}
\begin{ejercicios}
\item En un examen de 80 preguntas, Luisa acertó 62. ¿Qué porcentaje falló? \sol{18 de 80: $22.5\,\%$.}
\item El $30\,\%$ de los empleados de una empresa son técnicos y el $40\,\%$ de los técnicos son mujeres. Si hay 36 técnicas, ¿cuántos empleados hay? \sol{$12\,\%$ del total son 36: total 300.}
\item La matrícula pasó de \$800\,000 a \$920\,000. ¿Qué porcentaje subió? \sol{$\tfrac{120}{800} = 15\,\%$.}
\end{ejercicios}
\piso{Como en el examen}{opción múltiple}
\begin{ejercicios}
\item Si el $25\,\%$ de un número es 30, el $60\,\%$ del mismo número es
\opciones{72}{50}{18}{90}
\sol{A. El número es 120 y su $60\,\%$, 72.}
\item Un indicador pasó del $8\,\%$ al $10\,\%$. ¿Cuál afirmación es correcta?
\opcionesl{Subió $2\,\%$ y también $2$ puntos.}{Subió 2 puntos porcentuales, es decir, un $25\,\%$.}{Subió un $20\,\%$.}{Subió $10\,\%$.}
\sol{B. $2 \div 8 = 25\,\%$.}
\item En un salón, el $60\,\%$ son hombres. Si se van 6 hombres, quedan iguales hombres y mujeres. ¿Cuántos estudiantes había?
\opciones{30}{24}{36}{60}
\sol{A. Hombres $0.6x$, mujeres $0.4x$; $0.6x - 6 = 0.4x$: $x = 30$.}
\end{ejercicios}

\bitacora

% ------------------------------------------------------------------
\leccion{pu-pr-5}{Aumentos, descuentos e IVA}{Juego del factor escondido: subir y bajar es multiplicar.}

\subsection{La idea: cada cambio es un factor}
Aumentar un $25\,\%$ es pasar de $100\,\%$ a $125\,\%$: multiplicar por $1.25$. Descontar un $20\,\%$ es quedar en el
$80\,\%$: multiplicar por $0.80$. Pensar en \clave{factores} tiene dos ventajas enormes:
\begin{enumerate}
  \item Los cambios sucesivos se \emph{multiplican}: subir $25\,\%$ y bajar $20\,\%$ es $\times 1.25 \times 0.8 = \times 1$.
  \item Deshacer un cambio es \emph{dividir} por el factor: si el precio con IVA del $19\,\%$ es 238\,000, el precio sin
  IVA es $238\,000 \div 1.19 = 200\,000$, no $238\,000 \cdot 0.81$.
\end{enumerate}

\figura{factor}{Poner el IVA es multiplicar; quitarlo es dividir por el mismo factor.}

\subsubsection{Interés simple y compuesto}
Con interés \clave{simple} del $r\,\%$ anual, cada año se gana lo mismo: $I = C \cdot r \cdot t$. Con interés
\clave{compuesto}, cada año se aplica el factor $(1 + r)$ al saldo: $C (1 + r)^t$. Al cabo de dos años al $10\,\%$,
el simple da $20\,\%$ y el compuesto $21\,\%$.

\begin{regla}
\begin{itemize}
  \item Aumento del $p\,\%$: $\times\left(1 + \tfrac{p}{100}\right)$. Descuento del $p\,\%$: $\times\left(1 - \tfrac{p}{100}\right)$.
  \item Cambios sucesivos: se multiplican los factores. Factor final $f$: cambio total $(f - 1)\cdot 100\,\%$.
  \item Precio original $=$ precio final $\div$ factor.
  \item Descuentos sucesivos del $a\,\%$ y $b\,\%$ equivalen a uno de $\left(a + b - \tfrac{ab}{100}\right)\%$.
\end{itemize}
\end{regla}

\subsection{Ejemplos resueltos}
\begin{ejemplo}[Dos descuentos]
Una tienda ofrece $30\,\%$ de descuento y, en caja, un $10\,\%$ adicional. ¿Cuál es el descuento total?
\solucion $0.7 \cdot 0.9 = 0.63$: se paga el $63\,\%$, el descuento es del $37\,\%$ (no del $40\,\%$).
\end{ejemplo}

\begin{ejemplo}[Precio sin IVA]
Un celular cuesta \$952\,000 con IVA del $19\,\%$. ¿Cuál es el valor del IVA?
\solucion Sin IVA: $952\,000 \div 1.19 = 800\,000$. El IVA es $152\,000$.
\end{ejemplo}

\begin{ejemplo}[¿Cuánto hay que subir para recuperar?]
Una acción cayó un $20\,\%$. ¿Cuánto debe subir para volver a su valor?
\solucion Quedó en $0.8$; para volver a 1 hay que multiplicar por $\tfrac{1}{0.8} = 1.25$: subir un $25\,\%$.
\end{ejemplo}

\begin{ejemplo}[Interés compuesto]
Se depositan \$1\,000\,000 al $10\,\%$ anual compuesto. ¿Cuánto hay al cabo de 3 años?
\solucion $1\,000\,000 \cdot 1.1^3 = 1\,000\,000 \cdot 1.331 = 1\,331\,000$.
\end{ejemplo}

\begin{trampa}
Sumar porcentajes sucesivos («$30\,\% + 10\,\% = 40\,\%$») o quitar el IVA restando el $19\,\%$ del precio final. En los
dos casos la base cambió.
\end{trampa}

\begin{atajo}
Un aumento seguido de un descuento del mismo porcentaje $p$ siempre deja una pérdida de $\tfrac{p^2}{100}\,\%$:
subir y bajar $10\,\%$ pierde el $1\,\%$; $20\,\%$ pierde el $4\,\%$.
\end{atajo}

\begin{reto}
El precio de un artículo sube un $25\,\%$ y luego baja un $20\,\%$. Comparado con el precio original, el precio final
\opcionesd{sube un $5\,\%$}{es el mismo}{baja un $5\,\%$}{baja un $4\,\%$}
\end{reto}

\begin{retoresuelto}
Factores: $1.25 \cdot 0.80 = 1$. El precio final es exactamente el original. \textbf{Opción B.} Quien suma $+25 - 20$
elige la A; quien usa el atajo de «subir y bajar el mismo porcentaje» con un porcentaje que no es el mismo elige la D.
\end{retoresuelto}

\taller
\piso{Calienta}{factores}
\begin{ejercicios}[2]
\item Factor de un aumento del $8\,\%$. \sol{1.08.}
\item Factor de un descuento del $35\,\%$. \sol{0.65.}
\item $\$50\,000$ con $12\,\%$ de recargo. \sol{\$56\,000.}
\item Precio sin IVA de \$119\,000. \sol{\$100\,000.}
\end{ejercicios}
\piso{Entrena}{problemas}
\begin{ejercicios}
\item Un salario sube $10\,\%$ un año y $10\,\%$ el siguiente. ¿Cuánto subió en total? \sol{$1.1^2 = 1.21$: $21\,\%$.}
\item Tras un descuento del $15\,\%$, un pantalón cuesta \$68\,000. ¿Cuánto costaba? \sol{$68\,000 \div 0.85 = 80\,000$.}
\item Un comerciante compra a \$40\,000 y quiere ganar el $25\,\%$ \emph{sobre el precio de venta}. ¿A cuánto debe vender? \sol{La ganancia es $\tfrac14$ de la venta: el costo es $\tfrac34$ de ella. Venta $= 40\,000 \div 0.75 \approx 53\,333$.}
\end{ejercicios}
\piso{Como en el examen}{opción múltiple}
\begin{ejercicios}
\item Un precio baja $20\,\%$ y luego sube $20\,\%$. El precio final respecto al inicial
\opcionesd{es igual}{baja $4\,\%$}{sube $4\,\%$}{baja $2\,\%$}
\sol{B. $0.8 \cdot 1.2 = 0.96$.}
\item Una población crece $10\,\%$ cada año. Si hoy hay 20\,000 habitantes, dentro de 2 años habrá
\opciones{24\,000}{24\,200}{22\,000}{24\,400}
\sol{B. $20\,000 \cdot 1.21$.}
\item Con un descuento del $25\,\%$ pagué \$45\,000. Si el descuento hubiera sido del $40\,\%$, habría pagado
\opciones{\$27\,000}{\$36\,000}{\$33\,750}{\$30\,000}
\sol{B. Precio original $60\,000$; el $60\,\%$ es $36\,000$.}
\end{ejercicios}

\bitacora

% ------------------------------------------------------------------
\evaluacion
\begin{ejercicios}
\item Se reparten 450 dulces entre dos niños en razón $4 : 5$. ¿Cuántos más recibe el segundo?
\opciones{50}{90}{45}{250}
\sol{A. Partes de 50: 200 y 250. (4.1)}
\item En un mapa a escala $1 : 40\,000$, dos puntos están a 7.5 cm. La distancia real es
\opciones{3 km}{30 km}{300 m}{0.3 km}
\sol{A. $7.5 \cdot 40\,000 = 300\,000$ cm. (4.1)}
\item Si $a : b = 3 : 4$ y $b : c = 6 : 7$, entonces $a : c$ es
\opciones{$9 : 14$}{$3 : 7$}{$18 : 21$}{$1 : 2$}
\sol{A. $a : b : c = 9 : 12 : 14$. (4.1)}
\item Si 15 obreros hacen una carretera en 24 días, ¿cuántos días tardan 20 obreros?
\opciones{32}{18}{20}{16}
\sol{B. $360 / 20$. (4.2)}
\item 4 máquinas fabrican 600 tornillos en 5 horas. ¿Cuántos fabrican 6 máquinas en 8 horas?
\opciones{1\,200}{1\,440}{960}{1\,800}
\sol{B. 30 tornillos por máquina-hora; $48 \cdot 30$. (4.2)}
\item Un grifo llena un tanque en 5 horas y otro en 20 horas. Juntos tardan
\opciones{4 h}{12.5 h}{3 h}{2.5 h}
\sol{A. $\tfrac{100}{25}$. (4.3)}
\item Un tren sale a 60 km/h; una hora después sale otro, por la misma vía y en el mismo sentido, a 80 km/h. Lo alcanza al cabo de
\opciones{1 h}{2 h}{3 h}{4 h}
\sol{C. Ventaja 60 km, descuenta 20 km/h. (4.3)}
\item ¿Qué porcentaje de 250 es 40?
\opciones{$16\,\%$}{$40\,\%$}{$62.5\,\%$}{$6.25\,\%$}
\sol{A. $40/250 = 0.16$. (4.4)}
\item El $40\,\%$ del $25\,\%$ de 800 es
\opciones{80}{200}{320}{100}
\sol{A. $0.1 \cdot 800$. (4.4)}
\item Un televisor cuesta \$1\,428\,000 con IVA del $19\,\%$. Su precio sin IVA es
\opciones{\$1\,200\,000}{\$1\,156\,680}{\$1\,240\,000}{\$1\,409\,000}
\sol{A. $1\,428\,000 \div 1.19$. (4.5)}
\item Dos descuentos sucesivos del $20\,\%$ y del $25\,\%$ equivalen a uno de
\opciones{$45\,\%$}{$40\,\%$}{$5\,\%$}{$60\,\%$}
\sol{B. $0.8 \cdot 0.75 = 0.6$. (4.5)}
\item Un artículo bajó de \$80\,000 a \$60\,000. Para volver al precio inicial debe subir
\opciones{$25\,\%$}{$20\,\%$}{$33.\overline{3}\,\%$}{$30\,\%$}
\sol{C. $\tfrac{20\,000}{60\,000}$. (4.5)}
\end{ejercicios}
